\documentclass[10pt,a4paper]{amsart}
\usepackage{amsmath,amssymb,amsthm,graphicx}
\usepackage{verbatim}
\newenvironment{reference}{\comment}{\endcomment}
%\newenvironment{reference}{}{}
\newenvironment{slogan}{\comment}{\endcomment}
\newenvironment{history}{\comment}{\endcomment}

% For commutative diagrams we use Xy-pic
\usepackage[all]{xy}

% We use 2cell for 2-commutative diagrams.
\xyoption{2cell}
% Theorem environments.
%
\theoremstyle{plain}
\newtheorem{theorem}[subsection]{Theorem}
\newtheorem{proposition}[subsection]{Proposition}
\newtheorem{lemma}[subsection]{Lemma}

\theoremstyle{definition}
\newtheorem{definition}[subsection]{Definition}
\newtheorem{example}[subsection]{Example}
\newtheorem{exercise}[subsection]{Exercise}
\newtheorem{situation}[subsection]{Situation}

\theoremstyle{remark}
\newtheorem{remark}[subsection]{Remark}
\newtheorem{remarks}[subsection]{Remarks}

\numberwithin{equation}{subsection}

% Macros
%
\def\lim{\mathop{\mathrm{lim}}\nolimits}
\def\colim{\mathop{\mathrm{colim}}\nolimits}
\def\Spec{\mathop{\mathrm{Spec}}}
\def\Hom{\mathop{\mathrm{Hom}}\nolimits}
\def\Ext{\mathop{\mathrm{Ext}}\nolimits}
\def\SheafHom{\mathop{\mathcal{H}\!\mathit{om}}\nolimits}
\def\SheafExt{\mathop{\mathcal{E}\!\mathit{xt}}\nolimits}
\def\Sch{\mathit{Sch}}
\def\Mor{\mathop{\mathrm{Mor}}\nolimits}
\def\Ob{\mathop{\mathrm{Ob}}\nolimits}
\def\Sh{\mathop{\mathit{Sh}}\nolimits}
\def\NL{\mathop{N\!L}\nolimits}
\def\CH{\mathop{\mathrm{CH}}\nolimits}
\def\proetale{{pro\text{-}\acute{e}tale}}
\def\etale{{\acute{e}tale}}
\def\QCoh{\mathit{QCoh}}
\def\Ker{\mathop{\mathrm{Ker}}}
\def\Im{\mathop{\mathrm{Im}}}
\def\Coker{\mathop{\mathrm{Coker}}}
\def\Coim{\mathop{\mathrm{Coim}}}

% Boxtimes
%
\DeclareMathSymbol{\boxtimes}{\mathbin}{AMSa}{"02}

%
% Macros for moduli stacks/spaces
%
\def\QCohstack{\mathcal{QC}\!\mathit{oh}}
\def\Cohstack{\mathcal{C}\!\mathit{oh}}
\def\Spacesstack{\mathcal{S}\!\mathit{paces}}
\def\Quotfunctor{\mathrm{Quot}}
\def\Hilbfunctor{\mathrm{Hilb}}
\def\Curvesstack{\mathcal{C}\!\mathit{urves}}
\def\Polarizedstack{\mathcal{P}\!\mathit{olarized}}
\def\Complexesstack{\mathcal{C}\!\mathit{omplexes}}
% \Pic is the operator that assigns to X its picard group, usage \Pic(X)
% \Picardstack_{X/B} denotes the Picard stack of X over B
% \Picardfunctor_{X/B} denotes the Picard functor of X over B
\def\Pic{\mathop{\mathrm{Pic}}\nolimits}
\def\Picardstack{\mathcal{P}\!\mathit{ic}}
\def\Picardfunctor{\mathrm{Pic}}
\def\Deformationcategory{\mathcal{D}\!\mathit{ef}}
\usepackage{cleveref}
\crefname{originallemma}{original lemma}{original lemmas}
\Crefname{originallemma}{Original lemma}{Original lemmas}
\makeatletter
\newlabel{morphisms-lemma-fpqc-quotient-topology}{{02JY}{source}{}{sourcepointer-1}{}}
\newlabel{morphisms-lemma-fpqc-quotient-topology@cref}{{[originallemma][0][]02JY}{[1][0][]source}{}{}{}}
\makeatother
\crefname{originallemma}{original lemma}{original lemmas}
\Crefname{originallemma}{Original lemma}{Original lemmas}
\makeatletter
\newlabel{spaces-properties-lemma-cover-by-union-affines}{{03FX}{source}{}{sourcepointer-2}{}}
\newlabel{spaces-properties-lemma-cover-by-union-affines@cref}{{[originallemma][0][]03FX}{[1][0][]source}{}{}{}}
\makeatother
\crefname{originallemma}{original lemma}{original lemmas}
\Crefname{originallemma}{Original lemma}{Original lemmas}
\makeatletter
\newlabel{spaces-properties-lemma-quasi-compact-affine-cover}{{03H6}{source}{}{sourcepointer-3}{}}
\newlabel{spaces-properties-lemma-quasi-compact-affine-cover@cref}{{[originallemma][0][]03H6}{[1][0][]source}{}{}{}}
\makeatother
\title{A flat quasi-compact surjective morphism of algebraic spaces is universally submersive}
\author{The Stacks Project Authors}
\date{stacks-project:spaces-morphisms.tex:local-snapshot-20261012; GNU FDL 1.2 or later}
\begin{document}
\maketitle
\setcounter{subsection}{1}
\noindent\textit{Editorial note.} All delivered mathematical text is the author's own, copied from the pinned Stacks Project chapter file spaces-morphisms.tex; no statement or proof step was rewritten, reordered or removed.\par\medskip
\noindent\textit{Editorial note.} The proof works entirely with the presentations U -> V -> Y; the topology on |Y| used at the end is the quotient topology defined in the chapter, not an extra input.\par\medskip
\noindent\textit{Editorial note.} The three references of the counted spans are recorded as non-load-bearing navigation pointers; the argument that consumes them stays inside the counted proof span.\par\medskip
\noindent\textit{Editorial note.} Preamble: the three native excerpt spans of sources/stacks-spaces-morphisms/preamble.tex (lines 8-13, 14-18 and 154-217). The chapter file itself carries no preamble text.\par\medskip
\noindent\textit{Allowed standard prerequisites (editorial).}\begin{itemize}
\item Covering an algebraic space by a disjoint union of affine schemes. Properties of Spaces, Lemma spaces-properties-lemma-cover-by-union-affines, used at proof lines 6133-6136. Produces the etale cover V -> Y whose points define the topology of |Y|.
\item Affine etale covers of quasi-compact algebraic spaces. Properties of Spaces, Lemma spaces-properties-lemma-quasi-compact-affine-cover, used at proof lines 6137-6141. Produces the affine cover U\_i of V\_i x\_Y X.
\item The scheme case: flat quasi-compact surjective morphisms are submersive. Morphisms, Lemma morphisms-lemma-fpqc-quotient-topology, used at proof lines 6158-6159. Is the input applied to the morphism U -> V of schemes.
\item Definition of the topology on the underlying set of an algebraic space. used at proof lines 6154-6157. Turns openness of the preimages into openness in |Y|.
\end{itemize}\medskip
\ifdefined\hypertarget\hypertarget{sourcepointer-1}{}\fi
\noindent\textit{Source reference.} Original lemma 02JY: focused-seeds/AGP/sources/stacks-snapshot-20261011/tags.txt, source lines 3353--3353. Navigation pointer into the Stacks Project; the standard fact it supplies is applied inside the counted proof span and the argument resting on it stays there.\par\medskip
\ifdefined\hypertarget\hypertarget{sourcepointer-2}{}\fi
\noindent\textit{Source reference.} Original lemma 03FX: focused-seeds/AGP/sources/stacks-snapshot-20261011/tags.txt, source lines 4508--4508. Navigation pointer into the Stacks Project; the standard fact it supplies is applied inside the counted proof span and the argument resting on it stays there.\par\medskip
\ifdefined\hypertarget\hypertarget{sourcepointer-3}{}\fi
\noindent\textit{Source reference.} Original lemma 03H6: focused-seeds/AGP/sources/stacks-snapshot-20261011/tags.txt, source lines 4555--4555. Navigation pointer into the Stacks Project; the standard fact it supplies is applied inside the counted proof span and the argument resting on it stays there.\par\medskip
\begin{lemma}
\label{lemma-fpqc-quotient-topology}
Let $S$ be a scheme.
Let $f : X \to Y$ be a flat, quasi-compact, surjective morphism of
algebraic spaces over $S$.
A subset $T \subset |Y|$ is open (resp.\ closed) if and only
$f^{-1}(|T|)$ is open (resp.\ closed) in $|X|$.
In other words $f$ is submersive, and in fact universally submersive.
\end{lemma}

\begin{proof}
Choose affine schemes $V_i$ and \'etale morphisms $V_i \to Y$ such that
$V = \coprod V_i \to Y$ is surjective, see
Properties of Spaces,
Lemma \ref{spaces-properties-lemma-cover-by-union-affines}.
For each $i$ the algebraic space $V_i \times_Y X$ is quasi-compact.
Hence we can find an affine scheme $U_i$ and a surjective \'etale morphism
$U_i \to V_i \times_Y X$, see
Properties of Spaces,
Lemma \ref{spaces-properties-lemma-quasi-compact-affine-cover}.
Then the composition $U_i \to V_i \times_Y X \to V_i$ is a surjective,
flat morphism of affines.
Of course then $U = \coprod U_i \to X$ is surjective and \'etale
and $U \to V \times_Y X$ is surjective. Moreover, the morphism $U \to V$ is the
disjoint union of the morphisms $U_i \to V_i$. Hence $U \to V$ is surjective,
quasi-compact and flat. Consider the diagram
$$
\xymatrix{
U \ar[r] \ar[d] & X \ar[d] \\
V \ar[r] & Y
}
$$
By definition of the topology on $|Y|$ the set $T$ is closed
(resp.\ open) if and only if $g^{-1}(T) \subset |V|$ is closed
(resp.\ open). The same holds for
$f^{-1}(T)$ and its inverse image in $|U|$.
Since $U \to V$ is quasi-compact, surjective,  and flat we win by
Morphisms, Lemma \ref{morphisms-lemma-fpqc-quotient-topology}.
\end{proof}

\end{document}
